% Content of §2 Related Work: section header is in the master file.

\subsection{GIScience foundations for NL-to-spatial-database}

Natural language access to spatial databases sits within the GIScience programme articulated by \citet{goodchild1992gis}: formalising geographic concepts so they can be computed rigorously. The relational semantics we ground against trace to the $4$- and $9$-intersection formalisms~\citep{egenhofer1991topological}, the end-user topological relations of \citet{clementini1993small}, OGC Simple Feature Access~\citep{ogcsfs}, and GeoSPARQL~\citep{ogcgeosparql}. PostGIS~\citep{postgis} operationalises these ideas in PostgreSQL through geometry/geography types, casts, metric and topological functions, coordinate transformations, and KNN operators. Our grounding rules expose these constraints as prompt-side evidence.

\subsection{Natural language interfaces to spatial databases}

Early geospatial question answering focused on small fact databases: \citet{zelle1996geoquery} used inductive logic programming for GeoQuery, and \citet{punjani2018geospatial} extended this to template-based question answering over linked geospatial data, highlighting spatial-predicate coverage and coordinate-system awareness. Recent LLM-era work falls into two groups. The first targets end-to-end NL-to-executable-spatial-SQL: NALSpatial~\citep{liu2025nalspatial} translates user queries into spatial SQL (implemented over the SECONDO spatial DBMS rather than PostGIS); Monkuu~\citep{yu2025monkuu} is an IJGIS-published LLM-powered NL interface with dynamic schema mapping; GeoSQL-Eval~\citep{hou2026geosql} contributes a large-scale LLM evaluation on PostGIS functions and NL2GeoSQL tasks. The second targets higher-level GeoAI agents: GeoCogent~\citep{houcoding2025} is an LLM-based agent for geospatial code generation; GeoAgent~\citep{lin2026geoagent} is a hierarchical multi-agent architecture. Surveys map the broader opportunities~\citep{mai2024foundation}.

Our contribution is narrower than a general GeoAI agent and different from a function-coverage benchmark. We study grounding as an intervention in end-to-end execution-graded NL-to-PostGIS-SQL, where the semantic layer encodes OGC predicates, SRIDs, geography casts, KNN operators, and read-only bounded-output constraints. Compared with NALSpatial and Monkuu, we add a 125-question execution benchmark with an explicit Robustness suite, question-level McNemar testing, and an 11-family grounding-effect panel. Compared with GeoSQL-Eval, CQ-125 is smaller but designed for paired intervention analysis rather than broad function inventory.

Table~\ref{tab:related-comp} summarises the distinction.

\begin{table}[t]\centering\footnotesize
\caption{Direct comparison with recent NL-to-spatial-database systems. NALSpatial is implemented over SECONDO, not PostGIS. $\checkmark$: supported / reported; $-$: not reported.}
\label{tab:related-comp}
\setlength{\tabcolsep}{3pt}
\begin{tabular*}{\textwidth}{@{\extracolsep{\fill}}lccccc@{}}
\toprule
System & E2E & Function & Safety & Paired gnd. & Released \\
\midrule
Our framework                          & $\checkmark$ & $-$           & $\checkmark$ & $\checkmark$ & $\checkmark$ \\
NALSpatial~\citep{liu2025nalspatial}    & $-$           & $-$           & $-$           & $-$           & $\checkmark$ \\
Monkuu~\citep{yu2025monkuu}             & $\checkmark$ & $-$           & $-$           & $-$           & $-$           \\
GeoSQL-Eval~\citep{hou2026geosql}       & $-$           & $\checkmark$ & $-$           & $-$           & $\checkmark$ \\
GeoCogent~\citep{houcoding2025}         & $-$           & $-$           & $-$           & $-$           & $-$           \\
GeoAgent~\citep{lin2026geoagent}        & $-$           & $-$           & $-$           & $-$           & $-$           \\
\bottomrule
\end{tabular*}
\end{table}

\subsection{Text-to-SQL with LLMs and semantic layers}

LLM text-to-SQL has progressed through Spider~\citep{yu2018spider,rajkumar2022evaluating}, BIRD~\citep{li2023bird}, Spider~2.0~\citep{lei2025spider2}, and BEAVER~\citep{chen2024beaver}. DIN-SQL~\citep{pourreza2023dinsql} decomposes prompting into schema linking, classification, generation, and self-correction; DAIL-SQL~\citep{gao2024dail} uses skeleton-similarity few-shot selection; C3~\citep{dong2023c3} studies calibrated zero-shot prompting; CSpider~\citep{min2019cspider} is an early multilingual benchmark. Semantic-layer tools such as MetricFlow~\citep{metricflow} encode entities, measures, dimensions, and metric graphs for analytical warehouses. We adapt that vocabulary with GIS annotations, and use BIRD as a supporting warehouse track of the bilingual framework.

% =============================================================================
